The fixed-replay benchmark evaluates R-RRA and R-SHQ on the same 76 held-out
historical routes and five seeds as the common allocation comparison. Both
baselines use the same corrected RoundRobin pointer semantics and the same
leakage-free RF-optimized branching artifact as the cross-method comparison.
R-RRA completed $380/380$ admitted route instances and R-SHQ completed
$379/380$. R-RRA reached a mean cycle time of $29.2187$ h, mean waiting time of
$27.2885$ min, Gini of $0.5327$ and weighted fairness of $0.0053$. R-SHQ reached a lower observed mean
cycle time of $18.1068$ h and lower observed Gini of $0.4738$, but its mean
waiting time was higher at $27.9555$ min. Its weighted fairness was also lower
at $0.0032$. Weighted fairness is interpreted in the same direction as Gini
(lower is more balanced after availability weighting), while horizon-normalized
occupation is descriptive rather than universally better when lower or higher.

The paired ShortestQueue--RoundRobin comparison over the five seeds reports
one fewer completed route for R-SHQ, a mean cycle-time difference of
$-40003.00$ s, a mean waiting-time difference of $40.02$ s, a mean Gini
difference of $-0.0589$ and a mean weighted-fairness difference of $-0.00203$.
These are descriptive paired differences over five seeds, not a claim of
statistical superiority. The result supports R-SHQ as the stronger reactive
efficiency--fairness baseline among completed routes for this workload, with
the caveat that R-RRA completed one additional fixed-route instance and that
R-SHQ observes only current availability and cumulative assignment counts, not
future arrivals, future calendar states or processing-time predictions.
